\documentclass[10pt]{article}
\providecommand{\FIGDIR}{fig}
\newcommand{\DIGESTDATE}{31 October 2026}
\input{preamble}

\begin{document}

% ==================================================================== MASTHEAD
\thispagestyle{empty}
\begin{tikzpicture}[remember picture, overlay]
  \fill[Deep] (current page.north west) rectangle
      ($(current page.north east)+(0,-67mm)$);
  \fill[Amber] ($(current page.north west)+(0,-67mm)$) rectangle
      ($(current page.north east)+(0,-68.4mm)$);
  \foreach \x/\y/\r in {22/12/0.9, 41/26/1.5, 63/9/0.7, 88/31/1.1, 112/17/1.9,
                        138/29/0.8, 159/13/1.3, 178/54/1.0, 196/21/1.6,
                        30/44/0.6, 74/57/1.2, 125/48/0.9, 168/6/0.8, 191/59/0.7,
                        52/18/1.0, 100/60/0.6, 148/20/1.1, 66/49/1.4, 205/45/0.8}
    \fill[Sky, opacity=0.30] ($(current page.north west)+(\x mm,-\y mm)$) circle (\r mm);
\end{tikzpicture}

\vspace*{4mm}
{\sffamily\fontsize{7.6}{9}\selectfont\color{Sky}%
 AUTOMATED DAILY LITERATURE DIGEST\;\textbullet\;PhD RESEARCH WATCH\par}
\vspace{2.2mm}
{\sffamily\bfseries\fontsize{27}{29}\selectfont\color{white}%
 Particulate Matter\par}
{\sffamily\fontsize{27}{29}\selectfont\color{Sky}%
 Monitoring \& Health Effects\par}

\vspace{5.0mm}
{\sffamily\fontsize{8.4}{10}\selectfont\color{white}%
 Issue 2026-10-31\;\;\textcolor{Amber}{\textbar}\;\;Entry window 31 October 2026 (full day)\;\;%
 \textcolor{Amber}{\textbar}\;\;10 records screened in scope\par}
\vspace{18mm}

\noindent
\begin{minipage}[t]{0.190\linewidth}\metric{10}{RECORDS IN SCOPE\\(36 EXCLUDED)}{Deep}\end{minipage}\hfill
\begin{minipage}[t]{0.190\linewidth}\metric{6}{OF 10 SUBTOPIC\\BANDS POPULATED}{Teal}\end{minipage}\hfill
\begin{minipage}[t]{0.190\linewidth}\metric{2}{EFFECT ESTIMATES\\WITH A CI}{Sky}\end{minipage}\hfill
\begin{minipage}[t]{0.190\linewidth}\metric{2}{TIER-A\\RECORDS}{Amber}\end{minipage}\hfill
\begin{minipage}[t]{0.190\linewidth}\metric{2}{RECORDS CARRIED ON\\METADATA ALONE}{Clay}\end{minipage}

\vspace{6mm}

\section*{\sffamily\bfseries\color{Deep}Signal of the day}
\vspace{-1.5mm}{\color{Amber}\rule{\linewidth}{1.1pt}}\vspace{2.5mm}

\begin{keybox}
{\sffamily\bfseries\small\color{Deep}Idaho, 32 AQS sites over 20 years: winter wood smoke and the clean
background fell, fire-season PM$_{2.5}$ rose 2.54\,$\mu$g\,m$^{-3}$ per decade, and the annual mean stayed flat.}\\[1.4mm]
\footnotesize
4,260 site-months (2005--2024), Theil--Sen trends, PCA and IDW. Regional mean 7.53\,$\mu$g\,m$^{-3}$ (sites 3.16--12.08);
an August--September wildfire peak and a November--February inversion/wood-smoke peak, with July--September carrying
36.6\,\% of annual mass. 32.7\,\% of site-years exceed the 9.0\,$\mu$g\,m$^{-3}$ annual standard.
\textbf{Where it is weak:} an unreviewed preprint; IDW across 32 sites in mountain terrain is not an exposure surface,
and the two ``regimes'' are PCA loadings, not apportioned sources.
\textbf{Why it matters here:} a flat annual mean hiding two diverging seasonal sources is the case for season-specific
sensor correction --- smoke and wood-smoke aerosol need different optical calibrations.\par
\end{keybox}

\vspace{3mm}

\section*{\sffamily\bfseries\color{Deep}What changed today}
\vspace{-1.5mm}{\color{Amber}\rule{\linewidth}{1.1pt}}\vspace{2.5mm}

\footnotesize
\begin{itemize}

\item \textbf{Stroke in a clean karst city.} Guiyang 2016--2021, 11,517 strokes: PM$_{2.5}$ RR 1.026 (1.005--1.048) and
PM$_{10}$ 1.015 (1.001--1.030) per 10\,$\mu$g\,m$^{-3}$ at the maximal lag; single-pollutant models only.

\item \textbf{Informal workplaces.} Agra footwear cluster, 378 worker-sessions: no session met the WHO PM$_{2.5}$
guideline, and air velocity (0.80 to 0.16\,m\,s$^{-1}$ from factory to cottage) explained 79\,\% of between-tier variance.

\item \textbf{Antibiotics in PM$_{2.5}$.} A 74-compound UPLC-MS/MS method (LOQ to 0.37\,pg\,m$^{-3}$) found 22 antibiotics
in 29 Beijing samples; roxithromycin in all of them.

\item \textbf{Particulate matter leads T2DM attributable burden across Asia} (GBD 2023) --- a modelled attribution, not
new evidence of causality.

\item \textbf{A thin Saturday.} 46 candidates after 18 cross-day duplicates went to the 30 October issue; 10 admitted, two on
metadata alone (Guanzhong humidity--PM, kitchen purifiers in 338 cities). Retry those abstracts first.

\end{itemize}

\vspace{2mm}

\begin{keybox}
{\sffamily\bfseries\footnotesize\color{Deep}Window continuity}\\[1.2mm]
{\fontsize{7.2}{8.8}\selectfont
The 29 October issue closed with \texttt{last\_entry\_date} at 29 October. This issue was built on the
\textbf{1 November 05:2x UTC (00:2x CDT) run}; by the 22:00 rule 31 October is complete and 1 November is \emph{not}
opened. The run backfills 30 and 31 October as separate issues; candidates surfacing on both days (18) were screened
once and assigned to 30 October. This issue covers \textbf{31 October in full} (Saturday). The W44 weekly (25--31
October) and the October monthly are filed under this date.\par}
\end{keybox}

\vspace{3mm}
\par\vskip 0pt plus 18\baselineskip\penalty-200\vskip 0pt plus -18\baselineskip
\section*{\sffamily\bfseries\color{Deep}Corpus analytics}
\vspace{-1.5mm}{\color{Amber}\rule{\linewidth}{1.1pt}}\vspace{2.5mm}

\noindent\begin{minipage}{\linewidth}\includegraphics[width=\linewidth]{\FIGDIR/f1_subtopics.png}
\figcap{Subtopic assignment of the 10 in-scope records across six bands. Respiratory, exposure, occupational
\& indoor and burden carry two each; cardiovascular (Guiyang stroke) and sensing (antibiotics in PM$_{2.5}$) one each.
Neuro, reproductive, mechanistic and other-clinical bands are empty.}\end{minipage}

\vspace{2mm}
\noindent\begin{minipage}{\linewidth}\includegraphics[width=\linewidth]{\FIGDIR/f2_design_endpoint.png}
\figcap{Left --- architecture: three modelling (Bayesian influenza, causal-ML asthma, GBD), two
experimental (analytical method, filtration membrane), two metadata-only, and one each acute time-series,
cross-sectional (occupational) and measurement record analysis. Right --- endpoints from stroke and influenza to
asthma, T2DM burden, workplace IAQ and filtration efficiency.}\end{minipage}

\vspace{2mm}
\noindent\begin{minipage}{\linewidth}\includegraphics[width=\linewidth]{\FIGDIR/f3_metric_geography.png}
\figcap{Left --- particle metric: PM$_{2.5}$-only 6, unspeciated 3, joint PM$_{2.5}$+PM$_{10}$ 1
(Guiyang). Right --- geography: China 4, global/multi-country or laboratory 3, North America 2 (Idaho, Colorado), South
Asia 1 (Agra).}\end{minipage}

\vspace{2mm}
\noindent\begin{minipage}{\linewidth}\includegraphics[width=\linewidth]{\FIGDIR/f6_lifecourse.png}
\figcap{Life-course windows: one childhood-inclusive record (county asthma, all ages), one working-age
(footwear workers) and one older-adult signal (stroke risk stronger at $\geq$65). No in-utero or adolescent record.}\end{minipage}

\vspace{2mm}
\noindent\begin{minipage}{\linewidth}\includegraphics[width=\linewidth]{\FIGDIR/f4_heatmap.png}
\figcap{Subtopic $\times$ study architecture. The health bands sit in acute and modelling cells; sensing is a
single analytical-method cell; exposure splits between record analysis and metadata-only.}\end{minipage}

\vspace{2mm}

\noindent\begin{minipage}{\linewidth}\includegraphics[width=\linewidth]{\FIGDIR/f5_forest.png}
\figcap{Ratio scale, two estimates from one study: stroke incidence per 10\,$\mu$g\,m$^{-3}$ PM$_{2.5}$ and PM$_{10}$ at
the maximal lag (Guiyang, 11,517 strokes, single-pollutant GAM). Same population, same outcome --- not independent and
not poolable.}\end{minipage}

\vspace{4mm}

\par\vskip 0pt plus 24\baselineskip\penalty-200\vskip 0pt plus -24\baselineskip
\section*{\sffamily\bfseries\color{Deep}Paper-level digest}
\vspace{-1.5mm}{\color{Amber}\rule{\linewidth}{1.1pt}}\vspace{2.5mm}

{\fontsize{7.4}{9}\selectfont\color{Slate}Ordered by cluster. Each entry gives the
methodological core and the finding that matters, not an abstract paraphrase. Records
marked \textit{metadata only} carry no recoverable abstract and assert no finding.\par}

\band{Cardiovascular \& metabolic\hfill 1 record}{Teal}

\paperentry{Zhang short-term air pollution stroke incidence Guiyang karst city 2026}
{Time-series analysis of the impact of short-term exposure to atmospheric pollutants on stroke risk among residents in Guiyang City, 2016-2021}
{Front Public Health}
{Guiyang 2016-2021, 11,517 incident strokes (5.25/day), quasi-Poisson GAM with single-day and cumulative lags, controlling trend, weather, day-of-week and holidays. Per 10 $\mu$g\,m$^{-3}$ at the maximal lag: PM$_{2.5}$ RR 1.026 (1.005-1.048), PM$_{10}$ 1.015 (1.001-1.030); SO$_2$ 1.075, NO$_2$ 1.057, CO 1.346 per mg/m3. Women and adults >=65 carried more of the risk; under-65s responded only to NO$_2$ and PM$_{10}$. Limitation: single-pollutant models with the lag chosen at the maximum inflate the PM coefficients, co-pollutant collinearity is unaddressed, and the stroke series does not split ischaemic from haemorrhagic. Relevance: moderate - a low-concentration karst basin is where fixed-site averaging hides terrain-driven gradients that a dense low-cost network would resolve.}
{10.3389/fpubh.2026.1777772}

\par\vskip 0pt plus 10\baselineskip\penalty-200\vskip 0pt plus -10\baselineskip
\band{Respiratory \& allergic\hfill 2 records}{Sky}

\paperentry{Wang climate air pollutants influenza 48 countries convergent cross mapping 2026}
{Delayed and nonlinear associations and predictive causal signals between climate factors, air pollutants, and influenza epidemics in 48 countries}
{Glob Health Res Policy}
{48 countries, 2022-2024; FluNet weekly positivity against national climate and pollutant series in hemisphere-specific Bayesian distributed-lag models, then convergent cross mapping (CCM) for 'predictive causal' drivers. Northern Hemisphere peak OR 1.74 (1.34-2.26) at weekly cumulative SO$_2$ 3100 $\mu$g\,m$^{-3}$, lag 1.6 weeks; Southern peak for organic carbon OR 1.72 (0.65-4.58), null. CCM ranked dust PM$_{2.5}$ the leading pollutant driver (19 countries). Limitation: national-level exposure against surveillance positivity is fully ecological; three post-pandemic seasons cannot separate pollution from NPI relaxation, and CCM detects shared dynamics, not causation. Relevance: low - surveillance-scale; dust PM$_{2.5}$ is a size fraction low-cost optical sensors undercount.}
{10.1016/j.ghrp.2026.04.002}

\paperentry{Li spatiotemporal causal machine learning PM$_{2.5}$ asthma Colorado counties 2026}
{A Spatiotemporal Causal Machine Learning Framework for Quantifying the Impact of PM$_{2.5}$ on Asthma at the Community Level}
{GeoHealth}
{Colorado county-level monthly asthma hospitalisations and ED visits, 2011-2019, modelled with a spatiotemporal representation-learning network plus a causal-inference layer that estimates counterfactual outcomes under hypothetical PM$_{2.5}$ reductions. Estimated effects were heterogeneous, larger in communities with compounded environmental and social stressors, with wider uncertainty in data-sparse counties. No effect size per $\mu$g\,m$^{-3}$ is given in the abstract. Limitation: county-month aggregation with monthly PM$_{2.5}$ is ecological; 'causal' here rests on unconfoundedness given learned covariates, which a neural model cannot verify, and sparse rural counties are exactly where exposure error peaks. Relevance: moderate - the model's uncertainty is concentrated where monitors are absent, an explicit argument for sensor densification in rural Colorado.}
{10.1029/2025gh001735}

\par\vskip 0pt plus 10\baselineskip\penalty-200\vskip 0pt plus -10\baselineskip
\band{Exposure assessment \& modelling\hfill 2 records}{Clay}

\paperentry{Anjum Idaho PM$_{2.5}$ 2005-2024 wildfire vs wood-smoke regimes 2026}
{Long-Term Spatiotemporal Variability of Fine Particulate Matter (PM$_{2.5}$) Across Idaho, USA, Using Statistical and Geospatial Tools}
{Research Square (preprint)}
{EPA AQS daily PM$_{2.5}$ at 32 Idaho sites, 2005-2024, aggregated to 4,260 site-months; Mann-Kendall/Theil-Sen trends, PCA and IDW. Mean 7.53 $\mu$g\,m$^{-3}$ (site range 3.16-12.08). Bimodal climatology: an August-September wildfire peak and a November-February wood-smoke/inversion peak; July-September carries 36.6\% of annual mass. Cold-season and background series declined while fire-season PM$_{2.5}$ rose 2.54 $\mu$g\,m$^{-3}$ per decade, holding the annual mean flat; 32.7\% of site-years exceed the 9.0 $\mu$g\,m$^{-3}$ standard. Limitation: unreviewed preprint; IDW over 32 sites in mountainous terrain is not an exposure surface, and the 'regimes' are PCA loadings, not source apportionment. Relevance: high - two diverging seasonal regimes in inversion valleys are precisely where sparse FRM coverage fails and a calibrated PurpleAir layer (smoke-specific correction) adds most.}
{10.21203/rs.3.rs-11373778/v1}

\paperentry{Bei air moisture particulate matter winter Guanzhong Basin 2026}
{Contributions of air moisture to particulate matter pollution during wintertime in the Guanzhong Basin, China}
{Atmos Pollut Res}
{\textsc{metadata only.} Title and authors only (Bei, Chen, Jiang, Tie et al.). No abstract was retrievable at harvest. The title reports an analysis of how air moisture contributes to wintertime particulate pollution in the Guanzhong Basin. Design, mechanism, metrics and magnitude are not available; nothing is inferred. Relevance: potentially high - humidity-driven aerosol water uptake is the dominant error term in optical low-cost PM sensing, so a basin-scale moisture budget bears directly on correction. Retry abstract first.}
{10.1016/j.apr.2026.103256}

\par\vskip 0pt plus 10\baselineskip\penalty-200\vskip 0pt plus -10\baselineskip
\band{Sensing, forecasting \& instrumentation\hfill 1 record}{Amber}

\paperentry{Bei UPLC-MS/MS 74 antibiotics in PM$_{2.5}$ Beijing 2026}
{Optimization and application of a sensitive UPLC-MS/MS method for the quantitative determination of antibiotics in PM$_{2.5}$}
{J Chromatogr A}
{Method for 74 antibiotics in PM$_{2.5}$: T3 column, formic-acid mobile phase, Oasis HLB extraction with Na2EDTA. Recoveries 69.7-131.7\%, RSD 1.0-18.1\% at three spike levels, quantification limits down to 0.37 pg/m3. On 29 Beijing 24-h samples, 22 compounds exceeded detection limits (max 121.38 pg/m3); roxithromycin was in every sample and oxytetracycline had the highest mean (14.48 pg/m3). Limitation: 29 samples from one city and season give occurrence, not exposure; recoveries near 70\% for some classes bias low, and the link to resistance spread is asserted, not tested. Relevance: low for sensing - a filter-bound chemical fraction no optical sensor sees, but co-located filter sampling at sensor nodes could support it.}
{10.1016/j.chroma.2026.467564}

\par\vskip 0pt plus 10\baselineskip\penalty-200\vskip 0pt plus -10\baselineskip
\band{Occupational \& indoor\hfill 2 records}{Amber}

\paperentry{Arif Agra footwear cluster ventilation IAQ tiers 2026}
{Ventilation, indoor air quality and thermal comfort in India's leather footwear manufacturing cluster: a tier-resolved field study}
{Research Square (preprint)}
{Agra footwear cluster, three enterprise tiers (formal factory to cottage unit), 378 worker-sessions over four seasons and 147 days, breathing-zone measurements paired with ASHRAE 55 ballots. All pollutants rose linearly from large to cottage scale as air velocity fell from 0.80 to 0.16 m/s; velocity explained 79\% of between-tier pollutant variance. No session anywhere met the WHO PM$_{2.5}$ guideline; 91.3\% exceeded the ACGIH formaldehyde TLV. Limitation: unreviewed preprint; PM instrument and calibration are not described in the abstract, and session means against an annual guideline overstate comparability. Relevance: moderate - informal-sector workplaces with no monitoring are an obvious deployment niche for wearable or wall-mounted low-cost PM nodes.}
{10.21203/rs.3.rs-10780874/v1}

\paperentry{Zargarian spent coffee grounds MOF filtration membrane PM$_{2.5}$ bioaerosol photothermal 2026}
{Spent Coffee Grounds-Derived Micro-Nano-Architectured Air Filtration Membranes for Multiwavelength Light-Triggered Bioaerosol Inactivation}
{Small}
{Single-step co-electrospinning of MOF-177/polyacrylonitrile nanofibres with electrosprayed ZIF-8-decorated spent coffee ground microparticles. At 4.5x the NIOSH respirator test face velocity the membrane captured 99.7\% of PM$_{2.5}$ at about half the pressure drop of a MOF-free control and removed 83.4\% of ammonia; photothermal heating >80 C under simulated sunlight or NIR restored microbial suppression after three days of office air. Limitation: one challenge aerosol and short loading; no size-resolved efficiency at the MPPS, no long-term loading curve, and MOF stability under humidity is not shown. Relevance: low - a filtration material, not a measurement; at most a candidate inlet pre-filter for bioaerosol-sensitive sensor housings.}
{10.1002/smll.76528}

\par\vskip 0pt plus 10\baselineskip\penalty-200\vskip 0pt plus -10\baselineskip
\band{Burden, policy \& mitigation\hfill 2 records}{Coral}

\paperentry{Lee type 2 diabetes burden Asia GBD 2023 particulate matter leading risk 2026}
{The evolving burden of type 2 diabetes mellitus in Asia, 1990-2023: A Global Burden of Disease Study 2023}
{Diabet Med}
{GBD 2023 secondary analysis, 34 Asian countries in five regions, 1990-2023. Age-standardised T2DM prevalence and DALY rates rose almost everywhere, most in South Asia (prevalence +66.5\%, DALYs +75.6\%); High-income Asia Pacific DALYs fell 2.4\% (UI crosses zero). Particulate matter pollution was the leading attributable risk factor in every Asian region, though low-fruit diet led in South Asia (157 DALYs/100,000). Limitation: attribution inherits the GBD PM$_{2.5}$-diabetes exposure-response curve and satellite-model exposure surfaces; it is not new evidence of causality, and the abstract gives no PM-attributable numbers. Relevance: low - burden framing; the exposure layer is exactly the satellite-fusion product low-cost networks are used to correct in Asia.}
{10.1111/dme.70421}

\paperentry{Lu kitchen air purifier health-economic outcomes 338 Chinese cities 2026}
{Heterogeneous health-economic outcomes of kitchen air purifier interventions: A 338-city study in urban China}
{Environ Int}
{\textsc{metadata only.} Title and authors only (Lu, Yao, Jiang, Fu et al.). No abstract was retrievable at harvest. The title reports a 338-city assessment of heterogeneous health-economic outcomes of kitchen air-purifier interventions in urban China. Exposure model, health endpoints, cost basis and effect sizes are not available; nothing is inferred. Relevance: moderate - cooking-generated PM is the largest indoor source a low-cost sensor sees, and any intervention-scale cost-effectiveness estimate depends on indoor concentration inputs such sensors could supply. Retry abstract first.}
{10.1016/j.envint.2026.110587}

\clearpage
\section*{\sffamily\bfseries\color{Deep}Corpus register}
\vspace{-1.5mm}{\color{Amber}\rule{\linewidth}{1.1pt}}\vspace{2.5mm}

\input{table_rows}

\vspace{2mm}

\begin{keybox}
{\sffamily\bfseries\footnotesize\color{Deep}Provenance and method}\\[1.2mm]
{\fontsize{7.2}{8.8}\selectfont
Entry window \textbf{31 October 2026}, single day (Saturday). Harvester legs: \textbf{PubMed} health
\textbf{6}, sensing \textbf{0}; \textbf{Europe PMC} \textbf{35} (5 AGRICOLA retro-index records dropped);
\textbf{Crossref by ISSN} \textbf{22}; \textbf{arXiv} 150 kB, no in-window entry; \textbf{OpenAlex} HTTP 429. 62 unique,
2 already carried, 60 fresh; PubMed connector (EDAT 30--31 Oct, health and sensing axes) 9 PMIDs bucketed to 31 Oct,
1 merged, 4 appended (64); \textbf{18 cross-day duplicates assigned to 30 October}, \textbf{46 screened}. Abstracts for
10 PubMed records recovered by efetch; Crossref recovered none for the Elsevier day-0 deposits. \textbf{Consensus}
calibration sweep: 10 hits, 0 new (Yan 2026 and Guo 2026 matched by DOI to earlier issues/rejections).
\textbf{10 admitted}, \textbf{36 rejected} with logged reasons. \textbf{Trial watch} 30 October--1 November: 0
registrations updated. \textbf{2 carried on metadata alone}. \textbf{Scholar Gateway} not called.
Abstracts remain the copyright of their publishers.\par}
\end{keybox}

\end{document}
